% ============================================================
\section{Resumen}
\label{sec:resumen}
% ============================================================

Este informe documenta el Entregable~1 del servicio de asistencia en temas de investigación del SENAMHI orientado a evaluar la influencia del cambio climático sobre El Niño-Oscilación del Sur (ENOS) en las regiones oceanográficas Niño~3.4 y Niño~1+2. En cumplimiento de la actividad~(a) de los términos de referencia, el entregable cubre la revisión de metodologías de tiempo de emergencia (TOE, por sus siglas en inglés), la recopilación de la temperatura superficial del mar (TSM) observada y la recopilación de las simulaciones históricas y futuras de los modelos CMIP6. No produce todavía resultados de desempeño, proyección ni TOE, que corresponden a los Entregables~2 a~4, pero deja construido, verificado y versionado el insumo del que esos entregables dependen.

El producto operativo es un \emph{pipeline} reproducible (repositorio \archivo{smn_toe}, con un único punto de entrada, \texttt{run.sh}) que descarga, homogeneiza, enmascara y valida la variable \texttt{tos} (TSM, tabla CMOR \texttt{Omon}) de los modelos CMIP6 que publican los experimentos \texttt{historical}, \texttt{ssp245} y \texttt{ssp585}, más \texttt{ssp370} como escenario intermedio adicional, junto con la referencia observacional ERSSTv5 de NOAA. El universo de partida son los 102 modelos CMIP6 que publican \texttt{tos}/\texttt{Omon} en la federación ESGF; 41 de ellos publican los cuatro experimentos y 40 completaron la descarga y aprobaron el control de calidad (160 combinaciones modelo-experimento, Sección~\ref{sec:resultados}). Los 62 restantes quedaron excluidos por no publicar alguno de los experimentos en ninguna fuente consultada (44) o por disponibilidad o descarga incompleta (18).

El \emph{pipeline} incorpora controles específicos frente a problemas que contaminarían el dato sin producir errores visibles. Cuatro provienen de los datos y servicios de origen: valores de relleno que el regrillado no reconoce como faltantes (FGOALS-f3-L), filtrados por valor físico; archivos que declaran grados Celsius pero contienen valores en Kelvin (GISS-E2-1-G), detectados por el valor numérico; búsquedas paginadas en ESGF que un corte de red puede truncar sin aviso, que se repiten completas; y archivos de dos variantes de grilla en un mismo modelo, que se filtran por grilla. Además, el inventario exige explícitamente los cuatro experimentos y el control de calidad se recalcula en cada corrida. El detalle de cada control está en el documento de anexos.

En paralelo, el informe desarrolla la revisión científica: la partición de incertidumbre de \citet{hawkins2012}, su adaptación por \citet{szabo2016}, la aplicación a CMIP6 y Sudamérica de \citet{bruno2023}, el precedente institucional de \citet{alvarez2024}, la caracterización del Niño costero de \citet{takahashi2019} y un estado del arte de otras metodologías de TOE. De esa revisión se deriva la elección de los dos métodos de TOE que se aplicarán en el Entregable~4: M1 \citep{szabo2016} y \emph{Time of Detection} (ToD, \citealp{gopika2024}). Las instrucciones de instalación y uso del \emph{pipeline}, así como la descripción detallada de cada paso, se entregan en un documento de anexos separado.

% ============================================================
\section{Alcance}
\label{sec:alcance}
% ============================================================

Los términos de referencia del \emph{Servicio de asistencia en temas de investigación para evaluar la influencia del cambio climático en los eventos de El Niño y sus impactos en el Perú} (en adelante, TdR) definen el primer entregable, con plazo máximo de 25 días calendario, como el informe de la actividad~(a) del numeral~6: ``Realizar la revisión de metodologías sobre tiempo de emergencia (TOE), recopilación de datos de la información de temperatura superficial del mar (TSM), y simulaciones históricas y futuras de los modelos climáticos globales CMIP6 para los escenarios SSP2-4.5 y SSP5-8.5, correspondientes a las regiones Niño~3.4 y Niño~1+2''.

Para operacionalizar este mandato, el equipo técnico lo descompone en cinco tareas: (i) revisar los antecedentes científicos (ENOS y Niño costero, partición de incertidumbre, aplicaciones previas en Sudamérica y otras metodologías de TOE); (ii) consolidar el conjunto de modelos CMIP6 que publican la variable requerida; (iii) recopilar la referencia observacional ERSSTv5; (iv) definir los dominios oceanográficos y los periodos de análisis; y (v) homogeneizar unidades, calendarios, resoluciones espaciales y máscaras océano-tierra entre todas las fuentes. La tarea~(i) es de revisión bibliográfica y se desarrolla en la Sección~\ref{sec:marco}; las tareas (ii) a (v) son computacionales y se resuelven con el \emph{pipeline} descrito en las Secciones~\ref{sec:datos} a~\ref{sec:scripts}.

El control de calidad de este entregable verifica que el dato esté completo y homogéneo, no que los modelos representen bien el clima observado: esa evaluación de desempeño es materia del Entregable~2.
